% FILE: content/03_model.tex

\section{Ontologische Struktur des Modells}
\label{sec:model}

Dieser Abschnitt präsentiert den formalen Kern der Ontologie der Kontinua (OC).
Er liefert die einheitliche Notation, Axiome, Definitionen und strukturellen
Nebenbedingungen, die für die nachfolgenden Kapitel über Abschluss,
Synthese und praktische Anwendung benötigt werden.

OC beschreibt Kontinua in mehreren wissenschaftlichen Domänen mittels einer
einheitlichen strukturellen Sprache. Diese Sprache wird aus Achsen,
Potentialen, Flüssen, Schwellen, Grenzen, Zyklen und einem Maß der
Kontinuität aufgebaut. Jedes Kontinuum ist außerdem in einem umgebenden
Meta-Raum eingebettet. Alle Definitionen in diesem Abschnitt sind auf der
Ebene des formalen Vokabulars domänenunabhängig. Spätere Kapitel zur
Prüfung sortieren die Aussagen nach Beweisstatus und empirischer
Stützung; dieser Abschnitt liefert die gemeinsame Modellsprache, auf der diese
Prüfungen aufbauen.

\subsection{Axiomatische Grundlage: Niveau \texorpdfstring{$K_0$}{K\_0}}

Niveau \(K_0\) ist ein rein strukturelles Substrat.
Es ist kein physischer Raum; ihm ist weder Zeit, Energie, Geometrie noch Dynamik
innewohnend.
Seine Rolle ist es, die minimalen Bedingungen bereitzustellen, unter denen ein
Kontinuum höherer Ordnung logisch möglich ist.

\paragraph{Spezifikation von \texorpdfstring{\(K_0\)}{K\_0}}
\(K_0\) wird durch ein Tripel festgelegt:
\[
    K_0 = (S,\Delta,\mathcal{C}),
\]
mit:
\begin{itemize}
    \item \(S\) ist eine Menge von Zuständen;
    \item \(\Delta : S \times S \to \mathbb{R}_{\ge 0}\) ist eine strukturelle
          Differenzfunktion;
    \item \(\mathcal{C}\) ist eine strukturelle Relation (oder Familie von
          Relationen), die Unterscheidbarkeit erhält.
\end{itemize}
Auf dieser Ebene gibt es keine Zeitvariable und keinen Evolutionsoperator.

\paragraph{Axiom 0.1 (Unterschied und Unterscheidbarkeit).}
Für alle \(s_1,s_2 \in S\):
\[
    \Delta(s_1,s_2) = 0 \;\Rightarrow\; s_1 = s_2.
\]
Nichtverschwindende strukturelle Unterschiede bilden die minimale Bedingung für
Unterscheidbarkeit.
Ein Kontinuum kann nicht existieren, wenn es keine zwei unterscheidbaren
Zustände gibt.

\paragraph{Axiom 0.2 (Nichttriviale Schwelle \texorpdfstring{$\Theta_0$}{Theta\_0}).}
Es gibt ein \(\varepsilon > 0\), so dass
\[
    \forall s_1,s_2\in S,\qquad
    s_1 \neq s_2 \Rightarrow \Delta(s_1,s_2) \ge \varepsilon.
\]
Der Wert \(\Theta_0 = \varepsilon\) ist eine minimale strukturelle Schwelle:
Unterschiede kleiner als \(\varepsilon\) werden auf der Ebene von \(K_0\) nicht
aufgelöst.

\paragraph{Axiom 0.3 (Logisches Substrat).}
\(K_0\) trägt weder eine Zeitachse noch einen Dynamikoperator.
Es entfaltet sich nicht selbst und erzeugt nicht aus sich heraus höhere Ebenen.
Es legt nur logische Bedingungen für Unterscheidbarkeit und das
Vorhandensein nichttrivialer Unterschiede fest.
Alle Dynamik sowie alle Begriffe von Energie, Geometrie und Kausalität
gehören zu den Ebenen \(K_1\) und höher.

Insbesondere kann kein Operator auf \(K_0\) die Dimension erhöhen:
Unterschiede auf Ebene \(K_0\) werden immer entlang der vorhandenen
strukturellen Achse der Unterscheidbarkeit ausgedrückt.
Die mit \(K_0\) verbundene Dimensionalität ist durch den einbettenden
Meta-Raum \(M_0\) festgelegt.
Hier bezeichnet \(M_0\) den minimalen Meta-Raum, in dem das rein
strukturelle Substrat zur späteren Entstehung von \(K_0\to K_1\) auswertbar ist;
es ist kein physischer Behälter oder eine dynamische Umgebung.

\subsection{Konstruktion von Niveau \texorpdfstring{$K_1$}{K\_1}}

Niveau \(K_1\) ist das einfachste echte Kontinuum: Es führt Zeit,
eine eindimensionale Achse und eine elementare geometrische Struktur ein.

\paragraph{Spezifikation von \texorpdfstring{\(K_1\)}{K\_1}}
Das Kontinuum \(K_1\) wird definiert als
\[
    K_1 = \big(X_1,\tau_1,\Omega_1,A_1,P_1(t),J_1(t),\Theta_1,\partial\Omega_1,C_1,k_1(t)\big),
\]
mit:
\begin{itemize}
    \item \(A_1\) ist eine einzelne Achse (eine eindimensionale Koordinate);
    \item \((X_1,\tau_1)\) ist ein topologischer Raum, gewonnen aus den
          strukturellen Daten von \(K_0\), typischerweise ein Intervall
          \(X_1 = (a,b)\);
    \item \(\Omega_1\) ist eine Menge zulässiger Konfigurationen auf \(X_1\)
          mit passenden Regularitätsbedingungen.
          In diesem zeitabhängigen analytischen Beispiel gilt
          \[
              \Omega_1 = C^0(I,H^1(X_1,V_1)) \cap C^1(I,L^2(X_1,V_1)),
          \]
          wobei \(I\) das Zeitintervall und \(V_1\) der im Beispiel genutzte
          Hilbertraum ist; bei endlich vielen Freiheitsgraden ist ein
          endlichdimensionaler Vektorraum der spezielle Fall.
    \item \(P_1(t)\) ist ein energieähnliches Potentialfunktional auf
          \(\Omega_1\);
    \item \(J_1(t)\) sind Flüsse, die aus der Variation von \(P_1(t)\) folgen;
    \item \(\Theta_1\) kodiert minimale Bedingungen für klassische Stabilität;
    \item \(\partial\Omega_1\) ist die Grenze, an der Stabilitätsschwellen
          saturieren;
    \item \(C_1\) ist die Menge der klassischen Zyklen (etwa periodische Bahnen);
    \item \(k_1(t)\) misst eindimensionale Kontinuität nach der untenstehenden
          allgemeinen Definition.
\end{itemize}

\paragraph{Übergang \texorpdfstring{$\Psi_{0\to 1}$}{Psi\_0\_to\_1}}
Der Übergang \(K_0 \to K_1\) wird durch einen Operator erzeugt:
\[
    \Psi_{0\to 1} : (S,\Delta,\mathcal{C}) \mapsto (X_1,\tau_1,A_1,\Omega_1,\Theta_1),
\]
welcher das Skelettgerüst von \(K_1\) konstruiert; das Energiefunktional,
die Flussfamilie, die Grenze, die Zyklen und der Skalar der Kontinuität werden
in den nachfolgenden, \(K_1\)-spezifischen Klauseln festgelegt.
Der Skelettaufbau:
\begin{itemize}
    \item identifiziert eine Familie strukturell kompatibler Zustände in \(S\);
    \item stattet diese mit einer Ordnung und Topologie \((X_1,\tau_1)\) aus;
    \item definiert die erste Achse \(A_1\);
    \item konstruiert den zulässigen Konfigurationsraum \(\Omega_1\);
    \item induziert klassische Schwellen \(\Theta_1\).
\end{itemize}
Dies ist die erste Instanz dimensionaler Emergenz: eine kontinuierliche Achse
erscheint, die innerhalb des rein strukturellen Substrats von \(K_0\) nicht
repräsentierbar ist.

\subsection{Allgemeine Definition eines Kontinuums}

Für jedes Niveau \(K\) in der Hierarchie wird ein Kontinuum als das Tupel
\[
    K = \big(\Omega(K), A(K), P(t), J(t), \Theta(K), \partial\Omega(K), C(K), k(K,t)\big),
\]
definiert, wobei:
\begin{itemize}
    \item \(\Omega(K)\) eine nichtleere Menge zulässiger Zustände ist;
    \item \(A(K) = \{A_1,\dots,A_n\}\) ist eine endliche Menge von Achsen
von inkompatiblen Unterschieden;
    \item \(P(t)\) der Vektor der Potentiale auf \(\Omega(K)\) ist;
    \item \(J(t)\) ist der aggregierte Potenzial-Flussvektor; die
          Zustandsaktualisierungen werden durch den Evolutionsoperator \(E\)
          behandelt;
    \item \(\Theta(K) = \{\Theta_i\}\) ist die Menge der Schwellenfunktionen;
    \item \(\partial\Omega(K)\) ist die Grenze der zulässigen Zustände;
    \item \(C(K)\) ist die Familie strukturell stabiler Zyklen;
    \item \(k(K,t)\) ist das Maß der Kontinuität.
\end{itemize}

Jedes Kontinuum ist in einen umgebenden Meta-Raum \(M\) eingebettet, sodass
\[
    \Omega(K)\subseteq\Omega(M), \qquad A(K)\subseteq A(M).
\]
Der Meta-Raum liefert zusätzliche zulässige Zustände und Achsen, die künftige
dimensionale Erweiterungen von \(K\) aufnehmen können.

\paragraph{Zustandsraum und Grenze}
Schwellen werden als Funktionen \(f_i : \overline{\Omega(K)} \to \mathbb{R}\) mit
\[
    f_i(s) \le 0 \quad\text{für alle } s \in \Omega(K),
\]
definiert; die Grenze ist dann
\[
    \partial\Omega(K) = \{ s \in \overline{\Omega(K)} \mid \exists\, i: f_i(s) = 0 \}.
\]
Diese Definition benötigt keine Metrik und gilt gleichwertig für geometrische,
topologische, energetische, informationsbezogene oder logische Kontinua.

Die Grenze ist dynamisch. Die folgende Notation gibt eine Aktualisierungsabbildung
wieder, keinen Ableitungsbegriff eines mengenwertigen Objekts:
\[
    \partial\Omega(K,t+dt) =
    R\big(\partial\Omega(K,t),P(t),J(t),\Theta(K)\big),
\]
die \(\Omega(K)\) während Geburt, Lebenszeit und Todesereignissen schrumpfen,
expandieren oder bifurkieren kann.

\subsection{Taxonomie der Schwellen}

Jedes Kontinuum besitzt eine strukturierte Menge von Schwellen \(\Theta(K)\),
geordnet nach den folgenden Typen:

\begin{itemize}
    \item \textbf{Existenzschwellen} \(\Theta_{\mathrm{exist}}\): Bedingungen,
          unter denen \(\Omega(K)\neq\emptyset\).
    \item \textbf{Stabilitätsschwellen} \(\Theta_{\mathrm{stab}}\): Bedingungen
          für beschränkte Dynamik und nicht-divergente Flüsse.
    \item \textbf{Kritische Schwellen} \(\Theta_{\mathrm{crit}}\):
          Hypersurfaces, an denen qualitative Veränderungen auftreten
          (Phasenübergänge, Bifurkationen).
    \item \textbf{Dimensionsschwellen} \(\Theta_{\mathrm{dim}}\): Bedingungen,
          unter denen neue Achsen und höherdimensionale Kontinua entstehen können.
    \item \textbf{Todes-Schwellen} \(\Theta_{\mathrm{death}}\): strukturelle
          Grenzen, bei denen keine zulässigen Zustände mehr existieren und
          \(\Omega(K)\) kollabiert.
\end{itemize}

Das vollständige Schwellenfeld eines Kontinuums ist daher eine Sammlung von
Ungleichungen:
\[
    f_{i}^{(\mathrm{exist})}(s) \le 0,\quad
    f_{i}^{(\mathrm{stab})}(s) \le 0,\quad
    f_{i}^{(\mathrm{crit})}(s) \le 0,\quad
    f_{i}^{(\mathrm{dim})}(s) \le 0,\quad
    f_{i}^{(\mathrm{death})}(s) \le 0,
\]
mit den jeweiligen Randkomponenten als Gleichungen.
\subsection{Potentiale, Flüsse und strukturelle Spannung}

Potentiale \(P(t)\) kodieren die interne Konfiguration von Beschränkungen und
triebenden Kräften innerhalb eines Kontinuums.
Sie können sich in Energieoberflächen, chemischen Konzentrationen,
Membrandegre, repräsentationalen oder informationsbezogenen Strukturen sowie
institutionellen und normativen Spannungen ausdrücken.

Der aggregierte Flussvektor \(J(t)\) beschreibt die Änderungsrate der
Potentiale:
\[
    \frac{dP}{dt} = J(t),
\]
mit \(J(t)\) zerlegt in strukturell unterschiedliche Komponenten:
\begin{itemize}
    \item \textbf{stützende Flüsse} \(J_{\mathrm{support}}\), die Zyklen
          stabilisieren und \(k(K,t)\) stabil halten;
    \item \textbf{kritische Flüsse} \(J_{\mathrm{critical}}\), die das System auf
den kritischen Schwellenbereich \(\Theta_{\mathrm{crit}}\) und
          \(\Theta_{\mathrm{dim}}\) zu- oder über diese Schwellen treiben;
    \item \textbf{destruktive Flüsse} \(J_{\mathrm{kill}}\), die das System über
          \(\Theta_{\mathrm{death}}\) treiben und \(k(K,t)\) reduzieren.
\end{itemize}

Strukturelle Spannung \(T(K,t)\) ist ein Funktional von Potentialen, Achsen und
Gradienten (schematisch \(T=T(P,A,\nabla P)\)).
Sie misst, wie stark die aktuelle Konfiguration die Schwellenlandschaft belastet.
Dimensionsübergänge treten auf, wenn \(T(K,t)\) den Wert
\(\Theta_{\mathrm{dim}}(K)\) überschreitet; Kollaps entsteht, wenn
Destruktivflüsse zusammen mit der Spannung das System über
\(\Theta_{\mathrm{death}}(K)\) treiben.

\subsection{Zyklen und Kontinuität}

Zyklen \(C(K)\) sind geschlossene Trajektorien in \(\Omega(K)\), die in
endlichem Abstand von der Grenze bleiben und Kontinuität bewahren.
Auf struktureller Ebene gilt:
\[
    L(C) = \int_C d\ell_K < \infty,\qquad
    S(C) = \min_{s \in C} \delta_{\partial\Omega}(s) > 0,
\]
wobei \(d\ell_K\) das gewählte strukturelle Längenelement für das realisierte
Kontinuum ist und \(\delta_{\partial\Omega}(s)\) die entsprechende
Abstandsfunktion zur Grenze.
Die oben stehende Grenzdefinition ist metrikkorrekt; diese Zyklusabschätzungen
fügen eine strukturelle Norm nur dann hinzu, wenn quantitativer Abstand oder
Länge gemessen werden.
Bei metrischen Realisierungen kann diese Norm eine induzierte Distanz sein.
Bei nichtmetrischen Realisierungen ist sie eine am Niveau zulässige
Trennungsfunktionalität, die genau auf der Grenze verschwindet und innerhalb
strukturell innerer Zustände positiv ist.

Wenn das Kontinuum entwickelt, wird der zulässige Bereich als zeitabhängig
verstanden:
\(\Omega(K,t)\) bezeichnet den Zustandssraum von \(K\) zum Zeitpunkt
\(t\), während \(\Omega(K)\) den zugrunde liegenden
niveau-relativen zulässigen Raum bezeichnet, wenn kein Zeitschnitt angegeben ist.
Wenn ältere Formeln \(\Omega(K(t))\) schreiben, ist diese Notation als dasselbe
Zeitabschnittsobjekt \(\Omega(K,t)\) zu lesen.
Das Modell übernimmt die folgende einheitliche Definition der Kontinuität:
\[
    k(K,t)
    = \chi_{\Omega}(s_K(t);K,t)
      S_{\mathrm{axes}}(K,t)
      S_{\mathrm{cycles}}(K,t)
      S_{\mathrm{flows}}(K,t)
      S_{\mathrm{coh}}(K,t),
\]
mit:
\begin{itemize}
    \item \(s_K(t)\) ist der aktuelle Zustand von \(K\) zum Zeitpunkt \(t\),
          und \(\chi_{\Omega}(s_K(t);K,t)\) ist der Existenzindikator:
          er ist \(1\), wenn \(s_K(t)\in\Omega(K,t)\), sonst \(0\);
    \item \(S_{\mathrm{axes}}(K,t)\) quantifiziert die effektive Achsen-Sättigung,
          beispielsweise das Verhältnis des effektiven Rangs der aktiven Achsen
a zum maximalen möglichen Rang dieses Niveaus;
    \item \(S_{\mathrm{cycles}}(K,t)\) misst die Stärke und Effizienz stabiler
          Zyklen, etwa als gewichtete Summe zyklischer Effizienzen,
          normalisiert auf Maximalwerte;
    \item \(S_{\mathrm{flows}}(K,t)\) erfasst Flussstabilität. Mit aggregierten
          stützenden und destruktiven Größen \(\Phi_{\mathrm{support}}\) und
          \(\Phi_{\mathrm{kill}}\) gilt lokal die folgende Regel. Der
          Nullfluss-Zweig gilt nur, wenn der geprüfte Zustand keinen aktiven
          Flusskanal besitzt:
          \[
              S_{\mathrm{flows}}(K,t)=
              \begin{cases}
              \dfrac{\Phi_{\mathrm{support}}(t)}
                    {\Phi_{\mathrm{support}}(t)+\Phi_{\mathrm{kill}}(t)},
                    & \Phi_{\mathrm{support}}(t)+\Phi_{\mathrm{kill}}(t)>0,\\[0.7em]
              1, & \Phi_{\mathrm{support}}(t)=\Phi_{\mathrm{kill}}(t)=0.
              \end{cases}
          \]
          Damit gilt bei Verschwinden der Stützung bei positivem destruktivem Fluss:
          \(S_{\mathrm{flows}}=0\);
    \item \(S_{\mathrm{coh}}(K,t)\in[0,1]\) kodiert globale
          strukturelle Kohärenz; diese kann die Konnektivität relevanter Graphen,
          die Kompatibilität von Schwellen und Flüssen sowie das Fehlen
          widersprüchlicher Restriktionen umfassen.
\end{itemize}
Alle Multiplikationsfaktoren in der obigen Definition werden im lokalen Modell auf
\([0,1]\) normiert, bevor das Produkt ausgewertet wird.
Der Operator zur Aktualisierung der Kontinuität \(U\), der später in
Abschnitt~\ref{sec:operators-full} entwickelt wird, überführt \(k(K,t)\) auf
den nächsten rechtmäßigen Zeitschritt anhand dieser Faktoren.

Per Konstruktion gilt \(0 \le k(K,t) \le 1\).
Positives \(k(K,t)\) ist der Lebendigkeitsindikator. Er bündelt vier
gleichzeitige Aussagen: Der aktuelle Zustand ist zulässig, die aktiven Achsen
bleiben gestützt, der Zyklusfaktor und der Flussstabilitätsfaktor sind beide
positiv, und der Kohärenzfaktor ist nicht kollabiert. Der volle Lebendigkeitsstatus
benötigt zusätzlich die in Abschnitt~\ref{sec:collapse-rebirth} genannten Klauseln
zu Existenz, Schwelle und Identität.
Der Tod entspricht \(k(K,t)\to 0\) zusammen mit dem Kollaps von
\(\Omega(K)\).

\subsection{Evolutionsoperator}

Die Entwicklung eines Kontinuums wird auf struktureller Ebene durch einen Operator
dargestellt:
\[
    E: K(t) \mapsto K(t+dt),
\]
oder ausgeschrieben:
\[
    E:\ (\Omega,A,P,J,\Theta,\partial\Omega,C,k)
    \longrightarrow (\Omega',A',P',J',\Theta',\partial\Omega',C',k').
\]
Hier bezeichnet \(A'\) die nach Evolutionsschritt neue Achsenmenge desselben
Kontinuums im erweiterten Tupel.
Dies ist ein unärer Evolutionsoperator auf einem Zeitschnitt von \(K\). Der aktive
Meta-Raum gehört zur Umgebung. Wenn er lokal benannt werden muss, erscheint er als
Index-Parameter und nicht als zweites Argument von \(E\), damit die
Operator-Arity konstant bleibt.

Das Modell fasst Flussaktualisierungen, Schwellenprüfungen, Zyklusstruktur,
Grenzbewegungen, Achsenerweiterung und Kontinuitätsaktualisierungen nicht in
einem einzigen Differentialsystem zusammen. Diese Komponenten werden als
gekoppelte Bestandteile eines eingeschränkten strukturellen Evolutionsrahmens mit
folgenden Anforderungen behandelt:

\begin{itemize}
    \item \textbf{Konsistenz:} wenn \(s(t) \in \Omega(K(t))\), dann gilt entweder
          \(s(t+dt) \in \Omega(K(t+dt))\) oder der Übergang ist mit einer
          Schwellenüberschreitung verbunden;
    \item \textbf{Schwellenkonformität:} Flüsse folgen der Ungleichungsstruktur
          in \(\Theta(K)\), außer bei explizitem Überschreiten kritischer
          oder Todes-Schwellen.
    \item \textbf{Monotonie der Dimension:} wenn \(A'\) eine neue Achse enthält,
die nicht als Kombination früherer Achsen darstellbar ist, dann gilt
          \(\dim(K(t+dt)) > \dim(K(t))\); die Dimension kann nicht sinken.
\end{itemize}

Die Dynamik läuft solange, bis \(\Omega(K(t)) \neq \emptyset\).
Ist \(\Omega(K(t^\ast)) = \emptyset\), ist das Kontinuum tot und \(E\) kann
nicht mehr sinnvoll auf ihm wirken.

\subsection{Einbettung in Meta-Räume}

Jedes Niveau \(K_x\) ist in einen Meta-Raum \(M_x\) eingebettet, der zusätzliche
a zulässige Zustände und Achsen bereitstellt.
Auf struktureller Ebene wird dies durch
\[
    \Omega(K_x(t)) \subseteq \Omega(M_x(t)), \qquad
    A(K_x) \subseteq A(M_x)
\]
gegeben.
Das universelle Evolutionsaxiom kann dann schematisch geschrieben werden als
\[
    K_x(t+dt) = E_{M_x(t)}\big(K_x(t)\big),
\]
mit der Zusatzbedingung, dass bei Versagen der Zeitschnitt-Einbettungsbedingung
\(\Omega(K_x(t))\not\subseteq\Omega(M_x(t))\), das Kontinuum \(K_x\)
aufhört zu existieren. Ist der Zeitschnitt klar, verbergen die verkürzten
Notationen \(K_x\) und \(M_x\) den Zeitaspekt \(t\).
Die dimensionale Weiterentwicklung von \(K_x\) ist nur möglich, wenn der Meta-Raum
eine Achse \(A_{\mathrm{new}}\in A(M_x)\setminus A(K_x)\) enthält und die
strukturelle Spannung die zugehörige Dimensionsschwelle übersteigt. Dies ist die
Achsen-Ursprungs-Bedingung weiter unten: Neue Achsen können nicht allein aus
\(K_x\) generiert werden.

\subsection{Geburt von Kontinua}

Das Auftreten eines neuen Kontinuums \(K_{x+1}\) aus \(K_x\) ist ein
schwelleninduzierter Phasenübergang.
Strukturell erfolgt die Geburt, wenn folgende Bedingungen erfüllt sind:

\begin{enumerate}
    \item Eine neue Klasse von Unterschieden erscheint, die nicht über die bestehenden
          Achsen \(A(K_x)\) beschrieben werden kann;
    \item Die mit diesen Unterschieden verbundene strukturelle Spannung erfüllt
          \[
              T(K_x,t) > \Theta_{\mathrm{dim}}(K_x).
          \]
    \item Der Einbettungsraum \(M_x\) enthält mindestens eine Achse, die die
          neuen Unterschiede tragen kann (d.h.\ \(A_{\mathrm{new}} \in A(M_x)
          \setminus A(K_x)\)).
    \item Es gibt eine nichtleere Menge zulässiger Zustände
          \(\Omega(K_{x+1})\), kompatibel mit der neuen Achse und den
          Schwellen.
\end{enumerate}

Der Operator der dimensionsbedingten Geburt
\[
    \Psi_{x\to x+1}: K_x \to K_{x+1}
\]
ist minimal und irreversibel: Jede nichttriviale Entstehung der neuen Achse
konstituiert das neue Kontinuum \(K_{x+1}\), und die Dimension von
\(K_{x+1}\) kann nicht auf die von \(K_x\) zurückgeführt werden, ohne
\(\Omega(K_{x+1})\) zu zerstören.
Das lokale Modell implementiert somit dimensionsbezogene Monotonie. Das
zugehörige Ergebnis über das Ausbleiben spontaner Erzeugung ist separat in
Abschnitt~\ref{sec:theorem-spontaneous-dimension-creation} dokumentiert.

\subsection{Leben von Kontinua}

Ein Kontinuum \(K\) gilt im Zeitintervall als lebendig, wenn
\[
    \Omega(K(t)) \neq \emptyset,\qquad
    k(K,t) > 0,\qquad
    C(K(t)) \neq \emptyset,
\]
und wenn die stützenden Flüsse \(J_{\mathrm{support}}\) die destruktiven
Flüsse \(J_{\mathrm{kill}}\) über den relevanten Zyklen integriert überwiegen.

Leben entspricht daher der beständigen Existenz von:
\begin{itemize}
    \item stabilen Zyklen im endlichen strukturellen Abstand zur Grenze;
    \item ausreichenden stützenden Flüssen zur Aufrechterhaltung dieser Zyklen;
    \item Kohärenz zwischen Potentialen, Achsen und Schwellen;
    \item kontrolliertem kritischem Verhalten ohne Überschreiten von
          \(\Theta_{\mathrm{death}}\).
\end{itemize}

Diese Bedingungen werden auf jeder Ebene \(K_x\) unterschiedlich interpretiert,
doch das strukturelle Muster ist vom Protocellniveau bis zu Institutionen gleich.

\subsection{Tod von Kontinua}

Ein Kontinuum \(K\) stirbt zum Zeitpunkt \(t^\ast\), wenn
\[
    \Omega(K(t^\ast)) = \emptyset,
\]
was \(\chi_{\Omega}(s_K(t^\ast);K,t^\ast)=0\) impliziert und damit
\(k(K,t^\ast)=0\).
Äquivalent gilt:
\begin{itemize}
    \item alle stabilen Zyklen verschwinden, \(C(K(t^\ast)) = \emptyset\);
    \item kein Zustand erfüllt die Schwellenungleichungen gleichzeitig;
    \item jeder Versuch einer dynamischen Fortsetzung verletzt mindestens eine
          Existenz-Schwelle \(\Theta_{\mathrm{exist}}\).
\end{itemize}

\paragraph{Irreversibilität des Todes}
Ist \(\Omega(K(t^\ast)) = \emptyset\), gibt es keinen strukturellen Operator
innerhalb desselben Niveaus, der ein nicht-leeres \(\Omega(K)\) rekonstruieren
kann.
Jegliche scheinbare Wiederbelebung entspricht der Geburt eines neuen
Kontinuums \(K'\) mit eigenen Schwellen und Zustandsraum, nicht der
Fortführung des ursprünglichen \(K\).
Daher ist der Tod strukturell irreversibel.

\subsection{Interaktion von Kontinua}

Zwei Kontinua \(K_a\) und \(K_b\) interagieren über den
Interaktionsoperator
\[
    E_{\mathrm{int}} : (K_a,K_b) \mapsto (K_a',K_b'),
\]
der auf ihre kombinierten Potentiale, Flüsse, Schwellen und Achsen wirkt.
Strukturell werden mehrere Regime unterschieden:

\begin{itemize}
    \item \textbf{Konkurrenz:} Flüsse des einen Kontinuums behindern die
          Stabilisierung der Zyklen im anderen; typischerweise können
          \(k_a(t)\) und \(k_b(t)\) nicht beide unbegrenzt wachsen.
    \item \textbf{Parasitismus:} Ein Kontinuum entzieht einem anderen
          stützende Flüsse und erhöht das eigene \(k\) auf Kosten des Wirts.
    \item \textbf{Symbiose:} stützende Flüsse sind gekoppelt, sodass beide
          Kontinua ihre Kontinuität erhöhen und ihren zulässigen Raum
          erweitern.
    \item \textbf{Fusion:} Achsen und Potentiale verschmelzen zu einem neuen
          Kontinuum \(K_{\mathrm{fusion}}\) mit vereinten Achsen und einem neuen
          Zustandsraum \(\Omega_{\mathrm{fusion}}\).
\end{itemize}

Interaktion kann selbst Dimensionsübergänge induzieren, wenn kombinierte
Unterschiede neue, bisher ungenutzte Achsen erzeugen und die Spannung über
\(\Theta_{\mathrm{dim}}\) treiben.

\subsection{Vertikale Hierarchie \texorpdfstring{$K_0$–$K_{12}$}{K0--K12}}

Innerhalb dieses formalen Schemas kann die OC-Hierarchie \(K_0,\dots,K_{12}\)
wie folgt zusammengefasst werden:

\begin{itemize}
    \item \(K_0\): rein strukturelles Substrat unterscheidbarer Zustände und
          Unterschiede, ohne Zeit, ohne Dynamik.
    \item \(K_1\): eindimensionales klassisches Kontinuum mit
          Energie-Funktionalen und Grundstabilitätsschwellen.
    \item \(K_2\): physikalische Kontinua einschließlich Feldern,
          Phasenübergängen, Perkolation, Berezinskii--Kosterlitz--Thouless-ähnlichen
          Übergängen und Massenstruktur in der Feldtheorie.
    \item \(K_3\): chemische Kontinua einschließlich Reaktionsnetzen,
          reflexiv-autokatalytischen und ernährungsinduzierten Strukturen,
          Konzentrationen und Umgebungsparametern.
    \item \(K_4\): protocelluläre Kontinua mit Membranen, osmotischen
          und Krümmungsschwellen, internen/externalen Gradienten und
          metabolischen Teilräumen.
    \item \(K_5\): frühe neuronale und bioelektrische Kontinua mit Ionenkanälen,
          Erregbarkeitsschwellen, gradientgetriebenen Flüssen und Proto-Spikes.
    \item \(K_6\): kognitive Kontinua mit repräsentationalen Achsen,
          Bindung, internen Modellen, Gedächtnisschwellen und
          Vorhersageschwellen.
    \item \(K_7\): soziale Kontinua mit Vertrauensschwellen, institutionellen
          Zyklen, Kommunikation und Koordinationsflüssen.
    \item \(K_8\): zivilisatorische Kontinua mit großskaligen
          Infrastrukturen, systemischen Schwellen und Langstreckenzyklen.
    \item \(K_9\): theoretische Kontinua aus Theorien, Paradigmen,
          Ontologien und Modellfamilien mit Konsistenz- und Kohärenzschwellen.
    \item \(K_{10}\): metatheoretische Kontinua, die Vergleich, Evaluation
          und Transformation über Klassen von Theorien regeln.
    \item \(K_{11}\): rekursive metatheoretische Kontinua über
          Operatorfamilien, metalogische Aktualisierung und Rekursion im
          Modellraum.
    \item \(K_{12}\): globale Kohärenz-Kontinua, die die Kompatibilität über
          Familien von \(K_{11}\)-Strukturen hinweg begrenzen.
\end{itemize}

Jedes Niveau übernimmt die allgemeine Kontinuumsstruktur und ergänzt neue
Achsen, Potentiale, Schwellen und Zyklen spezifisch für seine Domäne.
Dieses Kapitel leitet nicht für jede domänenspezifische Repräsentationstheorie
jede Einzelheit im vollen Detail her; es dokumentiert das gemeinsame
strukturelle Rückgrat, das diese Instantiierungen respektieren müssen.

\subsection{Zusammenfassung}

Dieser Abschnitt hat das ontologische Rückgrat von OC präsentiert: die Axiome von
\(K_0\), die Konstruktion von \(K_1\), die allgemeine Kontinuumsdefinition,
die Taxonomie der Schwellen, die Rollen von Potentialen, Flüssen und
struktureller Spannung sowie die Definition von Zyklen und Kontinuität.
Er hat zudem den Evolutionsoperator, die Einbettung in Meta-Räume,
die strukturellen Bedingungen für Geburt, Leben und Tod, den Interaktionsoperator
und die vertikale Hierarchie \(K_0\)–\(K_{12}\) festgelegt.
Alle weiteren Entwicklungen in Core und Erweiterungsarbeiten stützen sich auf
diese Struktur als gemeinsame Grundlage.

